% Options for packages loaded elsewhere
\PassOptionsToPackage{unicode}{hyperref}
\PassOptionsToPackage{hyphens}{url}
\PassOptionsToPackage{dvipsnames,svgnames,x11names}{xcolor}
%
\documentclass[
  11pt,
  letterpaper,
]{article}
\usepackage{amsmath,amssymb}
\usepackage{setspace}
\usepackage{iftex}
\ifPDFTeX
  \usepackage[T1]{fontenc}
  \usepackage[utf8]{inputenc}
  \usepackage{textcomp} % provide euro and other symbols
\else % if luatex or xetex
  \usepackage{unicode-math} % this also loads fontspec
  \defaultfontfeatures{Scale=MatchLowercase}
  \defaultfontfeatures[\rmfamily]{Ligatures=TeX,Scale=1}
\fi
\usepackage{lmodern}
\ifPDFTeX\else
  % xetex/luatex font selection
  \setmainfont[]{Liberation Serif}
  \setmonofont[]{DejaVu Sans Mono}
\fi
% Use upquote if available, for straight quotes in verbatim environments
\IfFileExists{upquote.sty}{\usepackage{upquote}}{}
\IfFileExists{microtype.sty}{% use microtype if available
  \usepackage[]{microtype}
  \UseMicrotypeSet[protrusion]{basicmath} % disable protrusion for tt fonts
}{}
\makeatletter
\@ifundefined{KOMAClassName}{% if non-KOMA class
  \IfFileExists{parskip.sty}{%
    \usepackage{parskip}
  }{% else
    \setlength{\parindent}{0pt}
    \setlength{\parskip}{6pt plus 2pt minus 1pt}}
}{% if KOMA class
  \KOMAoptions{parskip=half}}
\makeatother
\usepackage{xcolor}
\usepackage[margin=1in]{geometry}
\usepackage{longtable,booktabs,array}
\usepackage{calc} % for calculating minipage widths
% Correct order of tables after \paragraph or \subparagraph
\usepackage{etoolbox}
\makeatletter
\patchcmd\longtable{\par}{\if@noskipsec\mbox{}\fi\par}{}{}
\makeatother
% Allow footnotes in longtable head/foot
\IfFileExists{footnotehyper.sty}{\usepackage{footnotehyper}}{\usepackage{footnote}}
\makesavenoteenv{longtable}
\setlength{\emergencystretch}{3em} % prevent overfull lines
\providecommand{\tightlist}{%
  \setlength{\itemsep}{0pt}\setlength{\parskip}{0pt}}
\setcounter{secnumdepth}{-\maxdimen} % remove section numbering
% definitions for citeproc citations
\NewDocumentCommand\citeproctext{}{}
\NewDocumentCommand\citeproc{mm}{%
  \begingroup\def\citeproctext{#2}\cite{#1}\endgroup}
\makeatletter
 % allow citations to break across lines
 \let\@cite@ofmt\@firstofone
 % avoid brackets around text for \cite:
 \def\@biblabel#1{}
 \def\@cite#1#2{{#1\if@tempswa , #2\fi}}
\makeatother
\newlength{\cslhangindent}
\setlength{\cslhangindent}{1.5em}
\newlength{\csllabelwidth}
\setlength{\csllabelwidth}{3em}
\newenvironment{CSLReferences}[2] % #1 hanging-indent, #2 entry-spacing
 {\begin{list}{}{%
  \setlength{\itemindent}{0pt}
  \setlength{\leftmargin}{0pt}
  \setlength{\parsep}{0pt}
  % turn on hanging indent if param 1 is 1
  \ifodd #1
   \setlength{\leftmargin}{\cslhangindent}
   \setlength{\itemindent}{-1\cslhangindent}
  \fi
  % set entry spacing
  \setlength{\itemsep}{#2\baselineskip}}}
 {\end{list}}
\usepackage{calc}
\newcommand{\CSLBlock}[1]{\hfill\break\parbox[t]{\linewidth}{\strut\ignorespaces#1\strut}}
\newcommand{\CSLLeftMargin}[1]{\parbox[t]{\csllabelwidth}{\strut#1\strut}}
\newcommand{\CSLRightInline}[1]{\parbox[t]{\linewidth - \csllabelwidth}{\strut#1\strut}}
\newcommand{\CSLIndent}[1]{\hspace{\cslhangindent}#1}
\ifLuaTeX
  \usepackage{selnolig}  % disable illegal ligatures
\fi
\usepackage{bookmark}
\IfFileExists{xurl.sty}{\usepackage{xurl}}{} % add URL line breaks if available
\urlstyle{same}
\hypersetup{
  pdftitle={What Constraints Cannot Decide},
  pdfauthor={Working manuscript --- author attribution pending},
  pdfkeywords={Italo Calvino, Oulipo, anticombinatorics, computational
creativity, narrative planning, fictional worlds, generative
AI, epistemic revision},
  colorlinks=true,
  linkcolor={black},
  filecolor={Maroon},
  citecolor={black},
  urlcolor={black},
  pdfcreator={LaTeX via pandoc}}

\title{What Constraints Cannot Decide}
\usepackage{etoolbox}
\makeatletter
\providecommand{\subtitle}[1]{% add subtitle to \maketitle
  \apptocmd{\@title}{\par {\large #1 \par}}{}{}
}
\makeatother
\subtitle{Calvino's Anticombinatorics and the Unfinished Ontology of
Fictional Worlds}
\author{Working manuscript --- author attribution pending}
\date{23 November 2026}

\usepackage{fancyhdr}
\usepackage{graphicx}
\pagestyle{fancy}
\fancyhf{}
\fancyfoot[L]{\fontsize{7}{9}\selectfont Calvino-Open-Worlds-2026-11-23 \textperiodcentered{} 000\_\_RETURN\_PATH.txt}
\fancyfoot[R]{\fontsize{8}{10}\selectfont\thepage\hspace{8pt}\includegraphics[width=0.19in]{MARK.pdf}}
\renewcommand{\headrulewidth}{0pt}
\setlength{\footskip}{28pt}
\begin{document}
\maketitle
\thispagestyle{fancy}
\begin{abstract}
Italo Calvino's aborted collaboration with Paul Braffort is often read
as a primitive anticipation of generative artificial intelligence: a
computer searches a vast space of possible stories, while a writer
supplies the decisive deviation. This article argues that the
experiment's more enduring problem is not the size of the search space
but the construction and governance of the representation being
searched. Calvino's surviving synopsis already specifies event-dependent
physical constraints, yet defers transformations of character; its rules
about secrets implicitly require identities and states absent from the
original character--action matrix; and its aesthetic constraints belong
neither to the fictional world nor to its characters. A parallel history
of AI storytelling---from TALE-SPIN to author-directed and multiagent
planning---shows that enlarging the ontology can represent motives,
beliefs, and changes without resolving aesthetic judgment. Contemporary
probabilistic generation does not simply replace anticombinatorics: it
relocates selective judgments into learned representations, human
feedback, prompts, and post-generation evaluation. The article proposes
a layered account of narrative construction that distinguishes immutable
observations, revisable world hypotheses, agent-relative beliefs, and
contestable aesthetic purposes. In fictional-world reconstruction, the
central question is therefore not which single world survives but which
commitments must remain stable and which should be permitted to change.
\end{abstract}

\setstretch{1.15}
\section{Introduction: a secret that is not in the
data}\label{introduction-a-secret-that-is-not-in-the-data}

A computer is instructed to reconstruct a story from four people
and twelve named actions. Each action must be assigned a pair of distinct characters. The
elementary representation looks manageable: twelve actions, twelve
available ordered pairs of distinct characters, and therefore
\(12^{12}\) possible assignments before any restrictions are introduced.
Yet the moment the machine is told that one character cannot obtain
\emph{the same secret} twice, its seemingly complete vocabulary fails.
The action has an actor and a patient; it has no field in which one
secret can be distinguished from another. The next problem is not how to
search faster. It is what kind of thing a secret must become for a
search to run at all.

This unnoticed representational demand is a productive way to return to
Italo Calvino's \emph{The Burning of the Abominable House} and his
collaboration with Paul Braffort. The familiar story of the experiment
emphasizes the vastness of the possibility space and the need for
constraints to eliminate impossible combinations. Calvino, on this
reading, anticipates generative AI by staging the progressive selection
of a particular work from innumerable alternatives. The computer would
handle routine combinations and leave the artist the consequential
\emph{clinamen}: the departure from a rule that turns an admissible
sequence into literature (Braffort 1997; Duncan 2012).

That comparison becomes misleading precisely where it seems easiest. A
combinatorial solver and a probabilistic language model do not perform
the same operation. The first needs an inventory of candidate
assignments and explicit compatibility procedures; the second learns
conditional distributions and generates sequences without enumerating
all complete stories. A diffusion model likewise learns reverse
transitions through noisy representations rather than rejecting entire
candidate images (Ho, Jain, and Abbeel 2020). The weak analogy---many
possibilities become one artifact---does not explain why some outputs
count as coherent worlds, why characters act, or why a reader would
care.

The stronger historical proposition is that anticombinatorics exposes a
problem of \textbf{representational commitments}. Every attempted
reduction of possibility depends on prior decisions about the entities
that exist, the state that persists, the observer whose knowledge
counts, and the assumptions that can be revised. The surviving
Calvino--Braffort material is unusually revealing because these
decisions come apart. It handles physical consequences, gestures toward
psychological transformation, requires semantic objects its first event
schema does not specify, and finally distinguishes the programmer's
aesthetic constraints from the constraints on fictional characters. The
projected novel did not appear, but these mismatches are evidence about
a research problem, not merely about inadequate processing power
(Braffort 1997).

This article follows that problem into an alternative history of
computational storytelling and returns to contemporary generative AI. It
argues for a distinction between the consistency of an imagined world
and the continuing freedom to revise one's hypothesis about that world.
The distinction matters especially when several images or episodes
appear to suggest a single place. Such evidence can warrant a shared
fictional world, several related worlds, or an irreducibly ambiguous
one. An algorithm that settles the issue in advance may destroy the very
relation the images were meant to make possible.

\section{I. The number and what it cannot
count}\label{i.-the-number-and-what-it-cannot-count}

Braffort's retrospective account separates objects often collapsed into
one: Calvino's published fiction, the synopsis given to Braffort, the
selection programs actually written, the 1977 presentation, and the
unfinished ambition of \emph{L'ordre dans le crime}. The paper presented
in 1977 was assembled from Calvino's own French translation of the story
and the working synopsis. Braffort reports that the experimental
programs introduced logical and semantic compatibility matrices but did
not tame the combinatorial explosion. A completed novel was not the
output of this machine (Braffort 1997). Distinguishing narrative
premise, specification, implementation, demonstration, and projection
prevents imagined capabilities from becoming historical achievements.

The elementary arithmetic makes the distinction tangible. Four named
people yield \(4\times3=12\) ordered nonreflexive actor--patient pairs.
Assign one such pair independently to each of twelve named actions, and
the number of assignments is \(12^{12}=8,916,100,448,256\). The number
printed in Calvino's story is different; Alessandrini identifies the
arithmetic discrepancy (Alessandrini 2023). The error does not undermine
the elementary combinatorics, but the correct number does not count
literary plots either. Two assignments may describe equivalent
narratives; one assignment may demand contradictory temporal states;
another may be formally permissible yet dramatically inert. Counting
assignments is not counting stories.

The synopsis already tries to close that gap. It limits repeated acts
and their reversals and demands that certain acts exclude later
possibilities. Most plainly, once a murder occurs, the victim can no
longer act or suffer another action in the sequence (Braffort 1997).
This is an event-dependent rule: whether an action is admissible depends
on what has previously happened. It is consequently inaccurate to say
that Calvino's original method knew only static constraints. It knew at
least one kind of changing state, the physical availability of a
character.

What the synopsis postpones is more precise. Characters might lose or
acquire subjective incompatibilities after actions they perform or
suffer. Calvino notes the possibility and explicitly declines, for the
moment, to explore it (Braffort 1997). The gap is not \emph{change
versus stasis} but one kind of change versus another. A murder updates a
discrete condition with a comparatively stable interpretation. Betrayal,
seduction, grief, or revelation may change what a character expects,
intends, fears, or can imagine doing, without offering an equally simple
transition rule. Nor need those changes coincide with outward acts. A
person might discover a secret without acting on it, or act on a false
assumption and never learn that it was false.

The action schema reveals still another omission. In the rules for
acquiring secrets, it matters whether the same secret has been obtained
again or whether a different secret is involved. But an event
represented only as (actor, action, patient) cannot express that
distinction. It needs at least an object identity and some state of
possession or knowledge. Once that object is added, further questions
arrive: whether a secret can be independently learned, falsely believed,
disclosed to several people, or cease to be secret. These are not
ornamental refinements. They determine which sequences the machine can
correctly exclude. Anticombinatorics first requires the invention of
things whose differences have to be judged.

\section{II. The missing state is not only
physical}\label{ii.-the-missing-state-is-not-only-physical}

A parallel computational history clarifies what this invention entails.
In 1977, the year of Calvino's Centre Pompidou presentation, James
Meehan published \emph{TALE-SPIN}. Instead of filtering a complete
inventory of crime assignments, the system generated stories through a
problem solver, an assertion mechanism, and an inference mechanism.
Agents pursue goals; their actions update descriptions of the physical
and interpersonal world; inferred consequences may produce new goals
(Meehan 1977). TALE-SPIN was not a completed answer to literary
creation. Its importance here is architectural. Narrative possibility
was organized as a sequence of world-state changes, not merely as a set
of survivors after elimination.

The chronology does not establish direct influence between Meehan and
Calvino or Braffort. It disproves a simpler genealogy in which explicit
combinatorics in the 1970s gives way only much later to the simulation
of dynamic worlds. A competing representational strategy existed at the
time. The distinction between the strategies is causal, not merely
historical. In a compatibility filter, a candidate combination may be
ruled out by a predicate. In a world simulator, events produce the very
state against which later predicates are evaluated. Filtering and
simulation can be combined, but they make different choices about what
must be given in advance (Meehan 1977; Braffort 1997).

In 1981 Natalie Dehn pressed a further objection. A writer does not
always establish a world and allow its inhabitants to generate a plot.
The writer may first want a confrontation, revelation, or reversal and
then invent a place, relationship, or prior event that makes the desired
scene possible. Author-directed generation therefore works backward as
well as forward: the narrative goal can induce the construction of
conditions under which it can occur (Dehn 1981). This possibility
defeats a common assumption shared by naive reconstruction and naive
simulation---that there must be one preexisting world whose hidden
details the process is recovering. In fictional creation, an apparently
missing cause may not be waiting to be found. It may have to be
invented.

Later planning research formalizes some of the omitted machinery. Riedl
and Young's IPOCL planner distinguishes causal coherence from character
intentionality. Actions must belong to causally intelligible plans and
to an explanation of why particular characters perform them (Riedl and
Young 2010). Brenner's continual multiagent planning extends the state
of a storyworld to perceptions, communication, commitments, beliefs,
goals, and revision. A character may know, suspect, or fail to know a
relevant precondition. Two agents can inhabit the same physical
environment and face different actionable possibilities because their
representations of that environment differ (Brenner 2010). These
developments sharpen the problem Calvino deferred: not merely whether
events are jointly possible, but for whom, at what time, and under which
beliefs.

Young's later HEADSPACEX proposal adds a particularly useful
distinction. A character may possess desires with incompatible ends,
even where the system limits the intentions that can be adopted as
commitments (Young 2017). Psychological contradiction cannot simply be
treated as logical failure. A character who wants to protect another
person and also reveal that person's secret is not necessarily
malformed. The coexistence of these desires may be the plot's central
pressure. A system that uses consistency checks indiscriminately may
rule out complex characters while faithfully enforcing its own type
system.

These advances do not make a simulated story aesthetically adequate.
Brenner expressly states that plot quality or interestingness is not an
explicit concept in the proposed system and that a model of author or
reader judgment would be required to pursue it (Brenner 2010). Better
world state, more precise belief state, and richer character agency do
not decide which completed sequence counts as a work of art. This is a
recurring separation, not a defect that disappears with each new layer
of representation.

\section{III. The programmer, the character, and the reader judge
different
things}\label{iii.-the-programmer-the-character-and-the-reader-judge-different-things}

Calvino's synopsis recognizes this separation. Its categories
distinguish objective restrictions, subjective constraints derived from
character personalities, and aesthetic constraints---also described as
subjective constraints of the programmer. It asks whether the latter two
can be considered together (Braffort 1997). The question is almost a
diagram of the unfinished project. A rule about whether a corpse can act
belongs to the state of the fictional world. A rule about whether a
character would betray a friend belongs to a model of that character. A
rule about whether the betrayal comes too early or creates a dull ending
belongs to the maker's aesthetic evaluation. These rules may overlap in
practice, but they do not derive their authority from the same place.

That difference matters for the \emph{clinamen}. Duncan shows that
Calvino's early account of a literary machine in ``Cybernetics and
Ghosts'' assigns a decisive role to reception, whereas the later
anticombinatorial account returns creative privilege to the authorial
deviation (Calvino 1986; Duncan 2012). Within Oulipian practice, a
stricter formulation of \emph{clinamen} requires more than an unusual
result: a compliant alternative must have been available and
deliberately rejected for an aesthetic purpose (Duncan 2012). This
condition cannot be met merely by selecting a low-probability token. The
sampled token is generally a valid result of the probabilistic sampling
process, not evidence that an agent considered a known constraint, found
a compliant solution, and chose to violate it.

The strict definition must itself be handled cautiously. Berkman argues
that Calvino's arrangements and revisions do not always instantiate
fixed procedures that mechanically generate the final prose. A diagram
may organize an evolving work as much as generate it, while the
\emph{clinamen} may be spatial, metaphorical, or present in the ongoing
revision of the arrangement (Berkman 2020). The alternative to equating
randomness with artistic freedom is therefore not to locate a single,
detectable exception in each work. It is to ask what rule was operative
for which actor at which stage, how it was changed, and what
counterfactual compliant alternative the evidence can actually support.

There is an additional temporal complication in contemporary AI. The
human labor of preference ranking, demonstration, filtering, and reward
modeling can alter the distribution of future model outputs before any
particular writer supplies a prompt (Ouyang et al. 2022). The evaluator
may not be the same person as the prompter, reviser, reader, or maker of
the training examples. The question of aesthetic agency becomes
distributed, but distribution does not make the distinctions disappear.
Rewarding instruction-following, preserving semantic consistency, and
judging literary consequence are different operations. An optimized
distribution can conceal, rather than remove, decisions about which
possibilities are considered desirable.

For this reason a simple labor reversal---then the human wrote, now the
machine writes---is insufficient. Calvino already considered a machine
that might change its own conventions, and later allowed the artist an
exceptional role. Contemporary generation likewise encompasses distinct
operations: pretraining, preference tuning, decoding, prompting,
evaluation, editing, publication, and reception. The significant
research task is to locate what each operation treats as fixed and what
it can still revise.

\section{IV. The world that evidence does not yet
determine}\label{iv.-the-world-that-evidence-does-not-yet-determine}

The same problem changes character when the input is not a sequence of
events but several images of an imagined place. Conventional
reconstruction asks how multiple observations correspond to one spatial
scene. Fictional images need not obey that presupposition. They may
depict different times, incompatible narrative realities, allegorical
versions of a place, or different locations joined by one name. More
pixels or more views do not decide which hypothesis ought to be adopted.
The question is not only the geometry implied by the image but the
ontology and revision policy of the fictional world.

Dehn's author-directed proposal supplies one alternative to the
assumption that a stable world must precede its episodes: the world may
be constructed retrospectively to support a desired action (Dehn 1981).
Another comes from the different functions of Calvino's \emph{Invisible
Cities}. Berkman's account of the novel's organization and
Venice-related interpretive pressure does not require the fifty-five
descriptions to resolve into one geometrically coherent Venice (Berkman
2020). A recurrence can be semantic or mnemonic, not merely spatial.
Inferring one map from every city description would be an interpretive
decision, not a neutral reconstruction of observations.

An especially provocative proposal is Cory Siler's open-world narrative
generation, which distinguishes facts already established to an audience
from further worlds it presently finds plausible (Siler 2022). Its key
idea is nonmonotonicity: evidence can defeat an assumption that
previously made a fictional explanation implausible. A dragon introduced
late in the story may not contradict any earlier observation; it may
overturn the reader's default belief that dragons do not exist in this
genre or setting. A purely subtractive possibility-space account cannot
capture the return of a possibility that was rejected only
provisionally. The distinction between an immutable observation and a
revisable default is therefore central. This is a proposal rather than a
demonstrated general solution, and its source remains a verification
priority.

The resulting model is neither indiscriminate openness nor the pursuit
of a single final answer. Consider three images that suggest one city
but disagree about a tower's location, the coastline, and a building's
interior. One interpretation treats them as views of a stable city from
different times; another treats them as distinct cities with a common
emblem; a third treats them as mutually inconsistent but narratively
related worlds. Each may preserve important aspects of the images.
Before reconstructing an environment, a system should distinguish what
was observed, what has been inferred, what has been stipulated for the
purposes of the story, and what remains negotiable.

This is also where the analogy with diffusion models should stop. A
diffusion model can generate vivid landscapes through learned reverse
dynamics (Ho, Jain, and Abbeel 2020). It does not follow from that fact
alone that a set of its images describes one causal, persistent,
navigable location. A scene may appear photographically resolved while
its spatial identity remains underdetermined. Conversely, a world with
relatively schematic images can possess a well-specified geography and
history. The aesthetic resolution of an image, the computational
convergence of a sampling trajectory, and the ontological determination
of a world are different kinds of closure.

\section{V. An architecture of qualified
commitments}\label{v.-an-architecture-of-qualified-commitments}

The combined evidence supports a layered model rather than one universal
consistency filter. The model is not attributed to Calvino or to any of
the later systems; it is an analytical proposal designed to expose
decisions usually hidden inside the language of narrative possibility.
It separates \textbf{observations} from \textbf{world hypotheses},
\textbf{world hypotheses} from \textbf{character beliefs},
\textbf{character beliefs} from \textbf{intentions}, and all of these
from \textbf{aesthetic criteria}. Aesthetic evaluation can propose
revisions, but those revisions must reveal whether they change a
previously observed fact, a defeasible inference, an agent's belief, or
a desired dramatic trajectory.

\begin{longtable}[]{@{}
  >{\raggedright\arraybackslash}p{(\columnwidth - 6\tabcolsep) * \real{0.2500}}
  >{\raggedright\arraybackslash}p{(\columnwidth - 6\tabcolsep) * \real{0.2500}}
  >{\raggedright\arraybackslash}p{(\columnwidth - 6\tabcolsep) * \real{0.2500}}
  >{\raggedright\arraybackslash}p{(\columnwidth - 6\tabcolsep) * \real{0.2500}}@{}}
\toprule\noalign{}
\begin{minipage}[b]{\linewidth}\raggedright
Layer
\end{minipage} & \begin{minipage}[b]{\linewidth}\raggedright
What it represents
\end{minipage} & \begin{minipage}[b]{\linewidth}\raggedright
What may change
\end{minipage} & \begin{minipage}[b]{\linewidth}\raggedright
What should not be silently changed
\end{minipage} \\
\midrule\noalign{}
\endhead
\bottomrule\noalign{}
\endlastfoot
Observed commitments & What the story has shown or explicitly stipulated
& New observations may be added & Earlier observations without disclosed
revision \\
World hypotheses & Possible causal, spatial, and social states &
Defaults, inferred details, alternate configurations & Their evidential
status \\
Character perspectives & What each agent believes, knows, desires,
intends & Beliefs, motives, goals, commitments & Attribution of one
character's belief to another \\
Aesthetic purposes & Desired rhythm, pattern, ambiguity, genre, surprise
& Preferences and priorities through revision & Their ownership and
history \\
\end{longtable}

The categories are not mutually independent in lived composition. An
aesthetic purpose may cause the author to invent a new event; the event
may alter what a character knows; that knowledge may activate a plan;
the plan may create an unforeseen outcome that changes the author's
initial purpose. The point of separation is not to make these processes
autonomous but to make their crossings accountable. What changes, who
authorizes the change, and which alternatives remain visible become
first-class questions.

The smallest computational test follows from the problem of the secret.
Begin with the same twelve action assignments and compare a
representation containing only actor, action, and patient with one that
additionally identifies secrets and tracks who knows them. Construct two
sequences identical at the first level but different in whether a second
acquisition concerns the same secret. If the primitive representation
cannot distinguish them, no amount of search-space pruning can apply the
stipulated rule correctly. The next test gives two characters different
beliefs about the same secret. The same world state now supports
different available plans. A third test introduces incompatible desires
and checks whether the system discards a potentially revealing character
conflict. A fourth introduces a new image whose apparent contradiction
can be resolved by revising a genre assumption, not an established fact.

These are proposed discriminating experiments, not reported results.
Their value is methodological. They locate an error's level before
turning it into a model-quality score. A semantic-object error cannot be
repaired by more sampling; an epistemic error cannot be repaired by
stricter global consistency; a genre-default error need not be repaired
by changing the fictional geography; and an aesthetic disagreement may
not be an error at all. The history of computational storytelling
suggests that different kinds of failure demand different
representational remedies.

\section{Discussion: anticombinatorics after
generation}\label{discussion-anticombinatorics-after-generation}

The strongest counterargument is that learned models might eventually
absorb all these distinctions statistically. A sufficiently trained
generative model can often maintain character identity, obey explicit
world constraints, construct plausible causal histories, and revise a
narrative in response to instructions. It may not need symbolic
representations of secrets, intentions, or observation commitments to
produce convincing outputs. This possibility should be taken seriously.
Explicit formalisms are not the only means of achieving a behavior, and
an architecture's internal representations need not resemble those used
to evaluate it.

Yet there is a difference between achieving an output in one trial and
exposing the conditions under which the output remains valid across
revisions, characters, and contradictory observations. If the user asks
why a once-impossible city can now contain a tower, the system must
either articulate the changed commitments or leave the explanation
opaque. The demand for an explicit account becomes especially strong in
interactive worlds, where characters and users will continue to act and
new material will be absorbed. Probabilistic fluency can postpone this
accounting. It cannot by itself determine whether the new tower
contradicts prior observation, replaces a provisional inference, or
deliberately introduces a rupture in the fiction.

Another counterargument is literary. Great literature may violate causal
logic, treat location as metaphor, or cultivate characters whose motives
never become clear. A perfectly explicit world model might therefore be
not only unnecessary but hostile to the work. That challenge is
consistent with the argument here. The proposed layers are an instrument
for preserving the right kinds of ambiguity, not a demand to resolve
everything into a database. A system can explicitly represent that two
incompatible versions remain live, or that a narrator's testimony is
disputed, without deciding which is true. The failure is not the absence
of a final answer; it is the unacknowledged conversion of an unresolved
question into an invisible fixed assumption.

Calvino's own suspicion of the comfort offered by finite combinatorial
order complicates any program of unlimited formalization (Calvino 1986).
The aim should not be to inventory every conceivable emotional, social,
or aesthetic variable. Instead, representations should emerge from the
operations a work requires, remain contestable, and be revised when a
new scene makes an earlier distinction inadequate. A model of a city may
begin with shared landmarks and later require competing histories. A
model of a character may begin with a goal and later require
ambivalence. A model of an object may begin as a prop and later become
the secret on which the plot turns. The research object is the
transformation of the representation itself.

\section{Limitations and unresolved
territory}\label{limitations-and-unresolved-territory}

This argument is a close interpretive and technical comparison, not a
reconstruction of Braffort's executable programs. The original source
code, exhaustive test outputs, and the complete working dossier were not
available in the materials used here. Braffort's account is primary
testimony written later, and claims about implementation must remain
narrower than claims about the synopsis (Braffort 1997). The related
planning systems were produced in different research communities;
resemblance and chronological proximity do not establish direct
influence. The discussion of Siler relies on the supplied zettel and a
source URL that was blocked during this preparation; its details should
be independently rechecked before publication.

The proposed layered model is not a complete computational architecture.
It lacks a general mechanism for forming or ranking world hypotheses, a
complete treatment of unreliable perception and deliberate deception,
and a non-arbitrary account of revising aesthetic criteria. It also
leaves unresolved the question of whether preserving contradictory
possible worlds improves an actual reader's experience. The paper
therefore proposes discriminating tests rather than claiming evidence of
aesthetic superiority.

\section{Conclusion: keep the possibility space
answerable}\label{conclusion-keep-the-possibility-space-answerable}

The lasting difficulty in Calvino's anticombinatorial project was not
simply that there were too many stories to search. The rules needed
entities the original representation could not identify, characters
whose states changed in ways the system deferred, and aesthetic
judgments that could not be read off the world's physical or
psychological constraints. Computational storytelling subsequently
developed richer representations of goals, beliefs, and revision, but
this progress repeatedly disclosed new questions of judgment rather than
dissolving them.

Contemporary generative AI changes the production of prose and images
without making these questions obsolete. It can generate a compelling
outcome without providing an explicit inventory of the possibilities
excluded or the assumptions fixed along the way. The decisive design
problem becomes how to preserve, contest, and revise those commitments
as the fiction continues. An imagined world is not necessarily a latent
object waiting for its last missing view. It can be a negotiated
construction whose incomplete identity makes further stories possible.
In that sense, the most fertile anticombinatorics may not be the
operation that narrows a vast field to one survivor. It may be the
discipline of knowing when a possibility was never represented, when it
was provisionally excluded, and when the work requires letting it
return.

\section{References}\label{references}

\begin{CSLReferences}{1}{0}
\bibitem[\citeproctext]{ref-alessandrini2023matematica}
Alessandrini, Paolo. 2023. {``La Matematica Di Italo Calvino.''}
\emph{Archimede}, no. 3: 157--63.
\url{https://riviste.mondadorieducation.it/archimede/wp-content/uploads/sites/2/2023/12/0060.ARCH3_23_ARch_Cultura_Alessandrini_157_163.pdf}.

\bibitem[\citeproctext]{ref-berkman2020clinamen}
Berkman, Natalie. 2020. {``Italo Calvino's Oulipian Clinamen.''}
\emph{MLN} 135 (1): 255--80.
\url{https://doi.org/10.1353/mln.2020.0001}.

\bibitem[\citeproctext]{ref-braffort1997ordre}
Braffort, Paul. 1997. {``L'ordre Dans Le Crime: Une Expérience
Cybernétique Avec Italo Calvino.''} \emph{Europe}, no. 815 (March).
\url{https://www.paulbraffort.net/litterature/critique/calvino_crime.html}.

\bibitem[\citeproctext]{ref-brenner2010creating}
Brenner, Michael. 2010. {``Creating Dynamic Story Plots with Continual
Multiagent Planning.''} In \emph{Proceedings of the Twenty-Fourth AAAI
Conference on Artificial Intelligence}, 1517--22. AAAI Press.
\url{https://cdn.aaai.org/ojs/7567/7567-13-11097-1-2-20201228.pdf}.

\bibitem[\citeproctext]{ref-calvino1986cybernetics}
Calvino, Italo. 1986. {``Cybernetics and Ghosts.''} In \emph{The Uses of
Literature: Essays}, translated by Patrick Creagh, 3--27. Harcourt Brace
Jovanovich.
\url{https://www.jfki.fu-berlin.de/academics/SummerSchool/Dateien2011/Reading_Assignments/iuli_reader2.pdf}.

\bibitem[\citeproctext]{ref-dehn1981story}
Dehn, Natalie. 1981. {``Story Generation After TALE-SPIN.''} In
\emph{Proceedings of the Seventh International Joint Conference on
Artificial Intelligence}, 16--17.
\url{https://ijcai.org/proceedings/1981/4}.

\bibitem[\citeproctext]{ref-duncan2012calvino}
Duncan, Dennis. 2012. {``Calvino, Llull, Lucretius: Two Models of
Literary Combinatorics.''} \emph{Comparative Literature} 64 (1).
\url{https://doi.org/10.1215/00104124-1539226}.

\bibitem[\citeproctext]{ref-ho2020denoising}
Ho, Jonathan, Ajay Jain, and Pieter Abbeel. 2020. {``Denoising Diffusion
Probabilistic Models.''} In \emph{Advances in Neural Information
Processing Systems}. \url{https://arxiv.org/abs/2006.11239}.

\bibitem[\citeproctext]{ref-meehan1977talespin}
Meehan, James R. 1977. {``TALE-SPIN, an Interactive Program That Writes
Stories.''} In \emph{Proceedings of the Fifth International Joint
Conference on Artificial Intelligence}, 91--97.
\url{https://ijcai.org/proceedings/1977/17}.

\bibitem[\citeproctext]{ref-ouyang2022training}
Ouyang, Long et al. 2022. {``Training Language Models to Follow
Instructions with Human Feedback.''} In \emph{Advances in Neural
Information Processing Systems}. Vol. 35.
\url{https://arxiv.org/abs/2203.02155}.

\bibitem[\citeproctext]{ref-riedl2010narrative}
Riedl, Mark O., and R. Michael Young. 2010. {``Narrative Planning:
Balancing Plot and Character.''} \emph{Journal of Artificial
Intelligence Research} 39: 217--67.
\url{https://doi.org/10.1613/JAIR.2989}.

\bibitem[\citeproctext]{ref-siler2022openworld}
Siler, Cory. 2022. {``Open-World Narrative Generation to Answer Players'
Questions.''} In \emph{Proceedings of the Eighteenth AAAI Conference on
Artificial Intelligence and Interactive Digital Entertainment}, 307--10.
AAAI Press.
\url{https://ojs.aaai.org/index.php/AIIDE/article/download/21981/21750/26033}.

\bibitem[\citeproctext]{ref-young2017sketching}
Young, R. Michael. 2017. {``Sketching a Generative Model of Intention
Management for Characters in Stories: Adding Intention Management to a
Belief-Driven Story Planning Algorithm.''} In \emph{Proceedings of the
AIIDE-17 Workshop on Intelligent Narrative Technologies}, 281--87. AAAI
Press.
\url{https://cdn.aaai.org/ojs/12986/12986-52-16503-1-2-20201228.pdf}.

\end{CSLReferences}

\section{Appendix A. Assembly
instrument}\label{appendix-a.-assembly-instrument}

The exact publication prompt supplied by the researcher, version
\textbf{15.55-AM}, is archived as
\href{_PROMPTS/assembly_slipcase_v15.55-AM.xml}{assembly instrument} and
embedded without redaction in \texttt{index.html}. Its SHA-256 appears
in \texttt{000\_\_PROMPTS.txt} and \texttt{\_SLIPCASE/MANIFEST.json}.
This pointer is used instead of printing the entire long prompt in the
scholarly manuscript. The paper's present claims are not evidence
supplied by that prompt.

\section{Appendix B. Making history}\label{appendix-b.-making-history}

\textbf{Provided:} The researcher supplied a November 2026 note about
Calvino and generative AI, twenty-four previously produced zettels, and
three successive POML research/assembly instructions in the visible
conversation. \textbf{Preserved:} Card payloads were staged, hashed, and
copied as root TXT and Markdown mirrors; their BibTeX variants and link
occurrences remain inspectable. \textbf{Retrieved:} During this
preparation, the Braffort HTML account, IJCAI's Meehan page, and the
Duncan and Brenner PDFs were inspected. \textbf{Derived:} The present
argument, graph, field arrangements, bibliographic normalization,
interface, and rebuilding code were composed by GPT-6.
\textbf{Verified:} Package integrity, local hashes, graph counts, and
PDF rendering are computational checks recorded separately.
\textbf{Unverified:} Source text fidelity against a machine-exported
conversation transcript, complete source coverage, and original project
code. No human review of the compiled publication is asserted.

\section{Appendix C. Replication
path}\label{appendix-c.-replication-path}

Open \texttt{index.html} from \texttt{file://} to inspect the cards and
their sources. Download the embedded reconstruction JSON. Verify payload
hashes against \texttt{ZETTELS.json}, and compare each root TXT card to
its \texttt{\_MD/} mirror. Use \texttt{\_SLIPCASE/RELATIONS.jsonl} to
trace each native address and
\texttt{\_SLIPCASE/BIBLIOGRAPHY\_VARIANTS.json} to inspect bibliography
discrepancies. The return identity is
\texttt{Calvino-Open-Worlds-2026-11-23}; a future merged archive must
treat SHA-256 payload identity, not original zettel ID or filename, as
the basis for deduplication. Source bodies marked \texttt{LINK\_ONLY}
were not bundled. The working paper remains a revisable interpretation,
not a new primary source.

\phantomsection\label{refs}


\end{document}
